\documentclass[11pt]{article}
\usepackage[margin=1in]{geometry}
\usepackage{amsmath,amssymb,amsthm}
\usepackage[T1]{fontenc}
\usepackage{lmodern}
\usepackage[hidelinks]{hyperref}
\newtheorem{theorem}{Theorem}
\title{A three-step chip-firing counterexample to Conjecture 00000001237}
\author{gaochengzhecpu}
\date{October 9, 2026}
\begin{document}
\maketitle
\begin{abstract}
On the directed triangle with three unit-length arcs, a single chip
travels around the cycle. Its least return period is three, whereas the
least common multiple of the arc lengths is one. The counterexample
has exactly one legal firing at each step, so it applies to both
parallel chip firing and one-vertex-at-a-time chip firing.
\end{abstract}

\section{The proposed formula}
The source defines the chip-firing period to be the length of the
recurrent oscillation and asserts that, on a directed cycle, this period
equals the least common multiple of the arc lengths.
We take the ordinary directed cycle on the labeled vertices
$V=\mathbb Z/3\mathbb Z$, with arcs
\[
 0\longrightarrow1,\qquad1\longrightarrow2,\qquad2\longrightarrow0.
\]
Every arc has length one. This is also a valid instance of the formula
if lengths are treated as explicit positive integral edge data.
The graph is strongly connected and each vertex has indegree and
outdegree one.

\section{The chip-firing dynamics}
A configuration is a function $c:V\to\mathbb N$.
A vertex $v$ can fire when $c(v)$ is at least its outdegree. Firing
subtracts that outdegree from $c(v)$ and sends one chip along each
outgoing arc. In our graph a firing moves one chip from $v$ to $v+1$.
Let $e_v$ denote the configuration with one chip at $v$ and zero elsewhere.

\begin{theorem}
Starting from $e_0$, chip firing has least positive return period three.
The least common multiple of the arc lengths is one. Hence the proposed
period formula is false.
\end{theorem}
\begin{proof}
In $e_v$, vertex $v$ is legal and both other vertices are illegal.
Its unique legal firing produces $e_{v+1}$. The entire orbit is therefore
\[
 (1,0,0)\ \longmapsto\ (0,1,0)\ \longmapsto\ (0,0,1)
 \ \longmapsto\ (1,0,0).
\]
The first three configurations are distinct. The first positive return
is thus at time three. Exactly one vertex is legal at every state in
this orbit, so simultaneous firing of all legal vertices yields the
same transitions as sequential firing.
On the other hand, the three arc lengths are $1,1,1$, whose least common
multiple is $1$. Since $3\neq1$, the asserted equality fails.
\end{proof}

\paragraph{Counting convention.}
The period here counts firings (equivalently, parallel time steps) until
the same configuration returns on the same fixed vertex set.
The configurations are not identified under graph rotations.
The source asks for the length of a recurrent oscillation, not for the
number of completed laps or the number of firings of one designated
vertex. Under the stated configuration-return notion the cycle above
has period three.

\section{Formal verification}
\texttt{lean/Main.lean} defines the actual finite arc set, its lengths,
outdegrees, configurations, legal vertices, and chip-firing updates.
The parallel update subtracts the chips sent by a legal vertex and sums
the chips received along incoming arcs; it is not defined by declaring
a three-cycle. A separate \texttt{fire} function implements the usual
single-vertex update.
\begingroup\raggedright
\begin{itemize}
\item \texttt{arc\_\allowbreak count}, \texttt{one\_\allowbreak outgoing\_\allowbreak arc}, and
      \texttt{one\_\allowbreak incoming\_\allowbreak arc} check the actual directed triangle.
      \texttt{directed\_\allowbreak cycle\_\allowbreak strongly\_\allowbreak connected} gives paths of
      length at most two between all vertex pairs.
\item \texttt{unique\_\allowbreak legal\_\allowbreak vertex} and
      \texttt{one\_\allowbreak chip\_\allowbreak transition} prove legality and the exact update.
      \texttt{parallel\_\allowbreak equals\_\allowbreak unique\_\allowbreak legal\_\allowbreak firing} relates the
      parallel and sequential conventions in this orbit.
\item \texttt{actual\_\allowbreak period} computes Mathlib's
      \texttt{Function.minimalPeriod} as three.
      \texttt{arc\_\allowbreak lengths\_\allowbreak lcm} computes the finite-set lcm as one.
      \texttt{conjectured\_\allowbreak period\_\allowbreak formula\_\allowbreak false} is their inequality.
\end{itemize}
\endgroup
Finite computations are checked by Lean's kernel through ordinary
\texttt{decide}. No auxiliary numerical script is required. No unproved
placeholders, custom axioms, or native evaluation are used as proofs.

\paragraph{Reproduction.}
In \texttt{lean/}, run:
\begin{verbatim}
lake build
lake env lean -DwarningAsError=true Main.lean
\end{verbatim}
Lean 4.19.0 and the exact Mathlib revision are pinned.
Run \texttt{tectonic main.tex} from the submission directory for the PDF.
The exact bilingual conjecture is in \texttt{SOURCE.md}.

\begingroup\raggedright
\begin{thebibliography}{9}
\bibitem{source} The Last Math Competition,
\href{https://github.com/The-Last-Math-Competition/The-Last-Math-Competition/blob/main/conjectures/00000001237.md}{Conjecture 00000001237}.
\bibitem{mathlib} The Mathlib Community, Lean 4 mathematical library,
\href{https://github.com/leanprover-community/mathlib4/commit/c44e0c8ee63ca166450922a373c7409c5d26b00b}{revision c44e0c8ee63c}.
\end{thebibliography}
\endgroup
\end{document}
